\setcounter{section}{29}

  \zwischenueberschrift{Die Krümmung auf einer riemannschen Mannigfaltigkeit}

Es sei $M$ eine
\definitionsverweis {riemannsche Mannigfaltigkeit}{}{}
und das
\definitionsverweis {Tangentialbündel}{}{}
$TM$ sei mit dem
\definitionsverweis {Levi-Civita-Zusammenhang}{}{}
versehen. Wir wollen den
\definitionsverweis {Krümmungsoperator}{}{}
$R$ in dieser Situation genauer verstehen. Er ist bei gegebenen Vektorfeldern $V,W$ eine Abbildung
\maabbdisp {R(V,W)} { C^2(M,TM) } { C^0(M,TM)
} {,}
d.h. er ergibt angewendet auf ein
\zusatzklammer {zweifach differenzierbares} {} {}
Vektorfeld $Z$ das Vektorfeld

\mavergleichskettedisp
{\vergleichskette
{ R(V,W)(Z)
}
{ =} { \nabla_V \nabla_W Z
}
{ } {
}
{ } {
}
{ } {
}
}
{}{}{.}
So gesehen handelt es sich insgesamt bei einer unendlich oft differenzierbaren Mannigfaltigkeit um  eine Abbildung
\maabbdisp {} {C^\infty(TM) \times C^\infty(TM) \times C^\infty(TM)  } { C^\infty(TM)
} {,}
wobei aber die Argumente nicht gleichberechtigt sind.

__NOEDITSECTION__

\inputfaktbeweis
{R^n/Offene Menge/Riemannsche Metrik/Levi-Civita-Zusammenhang/Krümmungsoperator/Berechnung/Fakt}
{Lemma}
{}
{

\faktsituation {Es sei

\mavergleichskette
{\vergleichskette
{ U
}
{ \subseteq }{ \R^n
}
{  }{
}
{    }{
}
{   }{
}
}
{}{}{}
\definitionsverweis {offen}{}{}
versehen mit einer
\definitionsverweis {riemannschen Metrik}{}{}
$g_{ij}$. Es sei $\nabla$ der
\definitionsverweis {Levi-Civita-Zusammenhang}{}{}
auf dem
\definitionsverweis {Tangentialbündel}{}{}

\mavergleichskette
{\vergleichskette
{ E
}
{ = }{ U \times \R^n
}
{  }{
}
{    }{
}
{   }{
}
}
{}{}{.}}
\faktfolgerung {Dann gilt für den
\definitionsverweis {Krümmungsoperator}{}{}

\mavergleichskettedisp
{\vergleichskette
{ R( \partial_i , \partial_j)( \partial_k)
}
{ =} { \sum_{\ell = 1}^n  R^\ell_{ijk} \partial_\ell
}
{ } {
}
{ } {
}
{ } {
}
}
{}{}{}
mit

\mavergleichskettedisp
{\vergleichskette
{ R^\ell_{ijk}
}
{ =} { {\partial_i } \Gamma^\ell_{jk} - {\partial_j } \Gamma^\ell_{ik}  +  \sum_{ a  = 1}^n  \Gamma^a_{jk} \Gamma_{i a}^\ell -\sum_{ a  = 1}^n   \Gamma^a_{ik}  \Gamma_{j a}^\ell
}
{ } {
}
{ } {
}
{ } {
}
}
{}{}{,}
wobei die

\mavergleichskettedisp
{\vergleichskette
{ \Gamma_{ij}^\ell
}
{ =} { { \frac{   1  }{  2 } } \sum_{\mu = 1}^n g^{\ell  \mu }(\partial_ig_{j\mu }+\partial_jg_{i\mu }-\partial_\mu g_{ij})
}
{ } {
}
{ } {
}
{ } {
}
}
{}{}{}
die
\definitionsverweis {Christoffelsymbole}{}{}
des Zusammenhangs bezeichnen.}
\faktzusatz {}
\faktzusatz {}

}
{

Dies ist ein Spezialfall von
[[Reelles Vektorbündel/Trivial/Linearer Zusammenhang/Krümmungsoperator/Beschreibung/Fakt|Lemma 28.2]].

}

__NOEDITSECTION__

\inputfaktbeweis
{Riemannsche Mannigfaltigkeit/Levi-Civita-Zusammenhang/Krümmungsoperator/Eigenschaften/Fakt}
{Lemma}
{}
{

\faktsituation {Es sei $M$ eine
\definitionsverweis {riemannsche Mannigfaltigkeit}{}{}
und das
\definitionsverweis {Tangentialbündel}{}{}
sei mit dem
\definitionsverweis {Levi-Civita-Zusammenhang}{}{}
versehen.}
\faktuebergang {Dann besitzt der
\definitionsverweis {Krümmungsoperator}{}{}
folgende Eigenschaften.}
\faktfolgerung {\aufzaehlungsechs{Es ist
\mathl{R(V,W)Z}{} linear in allen drei Komponenten.
}{Für

\mavergleichskette
{\vergleichskette
{ P
}
{ \in }{ M
}
{  }{
}
{    }{
}
{   }{
}
}
{}{}{}
hängt
\mathl{(R(V,W)Z) (P)}{} nur von
\mathl{V(P),W(P),Z(P)}{} ab.
}{Es ist

\mavergleichskettedisp
{\vergleichskette
{ R(V,W)
}
{ =} { -R(W,V)
}
{ } {
}
{ } {
}
{ } {
}
}
{}{}{}
}{Es ist

\mavergleichskettedisp
{\vergleichskette
{ R(V,W)Z+ R(W,Z)V+R(Z,V)W
}
{ =} { 0
}
{ } {
}
{ } {
}
{ } {
}
}
{}{}{.}
}{Es ist

\mavergleichskettedisp
{\vergleichskette
{ \left\langle R(V,W) T  ,  Z   \right\rangle
}
{ =} { - \left\langle R(V,W) Z  ,  T   \right\rangle
}
{ } {
}
{ } {
}
{ } {
}
}
{}{}{.}
}{Es ist

\mavergleichskettedisp
{\vergleichskette
{ \left\langle R(V,W) T  ,  Z   \right\rangle
}
{ =} { \left\langle R(T,Z) V  ,  W   \right\rangle
}
{ } {
}
{ } {
}
{ } {
}
}
{}{}{.}
}}
\faktzusatz {}
\faktzusatz {}

}
{

\aufzaehlungsechs{Die Linearität in $Z$ beruht darauf, dass der Zusammenhang linear ist. Die Linearität in $V$ und in $W$ beruht auf
[[Mannigfaltigkeit/Vektorbündel/R/Zusammenhang/Ableitung bezüglich Vektorfeld/Eigenschaften/Fakt|Lemma 24.10]]
und auf
[[Mannigfaltigkeit/Vektorfelder/Lie-Klammer/Eigenschaften/Fakt|Lemma 11.8]].
}{Für die Abhängigkeit in $V$ und $W$ folgt die Aussage aus
[[Reelles Vektorbündel/Linearer Zusammenhang/Krümmungsoperator/Eigenschaften/Fakt|Lemma 28.4]].
Um zu zeigen, dass auch die Abhängigkeit in $Z$ nur von $Z(P)$ abhängt, können wir von der lokalen Situation auf

\mavergleichskette
{\vergleichskette
{ U
}
{ \subseteq }{ \R^n
}
{  }{
}
{    }{
}
{   }{
}
}
{}{}{}
ausgehen und

\mavergleichskette
{\vergleichskette
{ V
}
{ = }{ \partial_i
}
{  }{
}
{    }{
}
{   }{
}
}
{}{}{,}

\mavergleichskette
{\vergleichskette
{ W
}
{ = }{ \partial_j
}
{  }{
}
{    }{
}
{   }{
}
}
{}{}{}
und

\mavergleichskette
{\vergleichskette
{ Z
}
{ = }{ h \partial_k
}
{  }{
}
{    }{
}
{   }{
}
}
{}{}{}
mit einer zweifach stetig differnezierbaren Funktion
\maabb {h} { U } {\R
} {}
ansetzen. Es ist dann nach
[[Differenzierbare Mannigfaltigkeit/R/Vektorbündel/Linearer Zusammenhang/Lokale Beschreibung/Fakt|Lemma 25.10]],
[[Mannigfaltigkeit/Vektorbündel/R/Linearer Zusammenhang/Leibnizregel/Fakt|Satz 25.5]]
und
dem [[Differenzierbarkeit/Satz von Schwarz/Fakt|Satz von Schwarz]]

\mavergleichskettealigndrucklinks
{\vergleichskettealigndrucklinks
{ R \left(  \partial_i, \partial_j  \right) \left(  h \partial_k  \right)
}
{ =} { \nabla_{\partial_i }  \left(  \nabla_{\partial_j }  \left(  h \partial_k  \right)  \right) - \nabla_{\partial_j } \left(  \nabla_{\partial_i }  \left(  h \partial_k  \right)  \right)
}
{ =} { \nabla_{\partial_i }  \left(  h \nabla_{\partial_j } \partial_k+ \left(  \partial_j h  \right)  \partial_k  \right) - \nabla_{\partial_j } \left( h \nabla_{\partial_i }  \partial_k  + \left(  \partial_i h  \right) \partial_k  \right)
}
{ =} { h \nabla_{\partial_i } \left(  \nabla_{\partial_j }  \partial_k   \right) +  \left(  \partial_i h  \right) \nabla_{\partial_j } \partial_k + \left(  \partial_j h  \right) \nabla_{\partial_i } \partial_k  +  \left(  \partial_i \partial_j h  \right) \partial_k
-h \nabla_{\partial_j } \left(  \nabla_{\partial_i } \partial_k   \right) -  \left(  \partial_j h  \right) \nabla_{\partial_i } \partial_k - \left(  \partial_i h  \right) \nabla_{\partial_j } \partial_k  -\left(  \partial_j  \partial_i h  \right) \partial_k
}
{ =} { h \nabla_{\partial_i } \left(  \nabla_{\partial_j }  \partial_k   \right)-h \nabla_{\partial_j } \left(  \nabla_{\partial_i } \partial_k   \right)
}
}
{}{}{,}
woraus hervorgeht, dass dies nur von $h(P)$ abhängt.
}{Ist klar aufgrund der Definition und wegen
[[Mannigfaltigkeit/Vektorfelder/Lie-Klammer/Eigenschaften/Fakt|Lemma 11.8]].
}{Aufgrund der
\definitionsverweis {Torsionsfreiheit}{}{}
\zusatzklammer {siehe
[[Riemannsche Mannigfaltigkeit/Levi-Civita-Zusammenhang/Eigenschaften/Fakt|Satz 26.14]]} {} {}
und
der [[Vektorfelder/Lie-Klammer/Jacobi-Identität/Fakt|Jacobi-Identität]]
ist

\mavergleichskettealigndrucklinks
{\vergleichskettealigndrucklinks
{ R(V,W)Z+ R(W,Z)V+R(Z,V)W
}
{ =} { \nabla_V  \left(  \nabla_W Z  \right) -\nabla_W  \left(  \nabla_V Z  \right) - \nabla_{[V,W]} Z +\nabla_W  \left(  \nabla_Z V  \right) -\nabla_Z  \left(  \nabla_W V  \right) - \nabla_{[W,Z]} V+\nabla_Z  \left(  \nabla_V W  \right) -\nabla_V  \left(  \nabla_Z W  \right) - \nabla_{[Z,V]} W
}
{ =} { \nabla_V  \left(  [W,Z ]  \right) + \nabla_W  \left(  [Z, V]  \right) + \nabla_Z  \left(  [V, W]  \right) - \nabla_{[V,W]} Z - \nabla_{[W,Z]} V  - \nabla_{[Z,V]} W
}
{ =} {\nabla_V  \left(  [W,Z ]  \right)  - \nabla_{[W,Z]} V + \nabla_W  \left(  [Z, V]  \right) - \nabla_{[Z,V]} W + \nabla_Z  \left(  [V, W]  \right) - \nabla_{[V,W]} Z
}
{ =} { [V, [W,Z]] + [W, [Z,V]] + [Z, [V,W]]
}
}
{
\vergleichskettefortsetzungalign
{ =} { 0
}
{ } {}
{ } {}
{ } {}
}{}{.}
}{Wir können und aus Vektorfelder $V,W$ mit

\mavergleichskette
{\vergleichskette
{ [V,W]
}
{ = }{ 0
}
{  }{
}
{    }{
}
{   }{
}
}
{}{}{}
beschränken. Es ist dann nach
[[Riemannsche Mannigfaltigkeit/Levi-Civita-Zusammenhang/Eigenschaften/Fakt|Satz 26.14  (2)]]
und
dem [[Differenzierbarkeit/Satz von Schwarz/Fakt|Satz von Schwarz]]

\mavergleichskettealigndrucklinks
{\vergleichskettealigndrucklinks
{ \left\langle R(V,W)T , Z  \right\rangle
}
{ =} { \left\langle  \nabla_V \nabla_WT -\nabla_W \nabla_VT   ,  Z  \right\rangle
}
{ =} { \left\langle  \nabla_V \nabla_WT , Z  \right\rangle - \left\langle  \nabla_W \nabla_VT  ,  Z  \right\rangle
}
{ =} { - \left\langle  \nabla_WT , \nabla_V Z  \right\rangle + D_V\left\langle  \nabla_W T  ,  Z  \right\rangle + \left\langle  \nabla_VT  , \nabla_W  Z  \right\rangle -  D_W \left\langle  \nabla_V T  ,  Z  \right\rangle
}
{ =} { \left\langle  T ,  \nabla_W  \nabla_V Z  \right\rangle - D_W \left\langle  T ,  \nabla_VZ   \right\rangle  + D_V \left(  - \left\langle  T  ,  \nabla_WZ  \right\rangle +  D_W \left\langle T  ,  Z   \right\rangle  \right)  - \left\langle  T  ,  \nabla_V \nabla_W  Z  \right\rangle +  D_V \left\langle  T  ,  \nabla_W Z  \right\rangle - D_W  \left(  - \left\langle T  ,  \nabla_V Z   \right\rangle + D_V \left\langle T  ,  Z   \right\rangle  \right)
}
}
{
\vergleichskettefortsetzungalign
{ =} { \left\langle  T ,  \nabla_W  \nabla_V  Z  \right\rangle - \left\langle  T  ,  \nabla_V \nabla_W  Z  \right\rangle
}
{ =} { - \left\langle R(V,W)Z , T  \right\rangle
}
{ } {}
{ } {}
}{}{.}
}{Unter Verwendung von (3), (4) und (5) ist

\mavergleichskettealigndrucklinks
{\vergleichskettealigndrucklinks
{ \left\langle R(V,W)T , Z  \right\rangle
}
{ =} { - \left\langle R(T,V)W , Z  \right\rangle - \left\langle R(W,T)V , Z  \right\rangle
}
{ =} { \left\langle R(T,V) Z , W  \right\rangle + \left\langle R(W,T) Z ,  V  \right\rangle
}
{ =} { - \left\langle R(V,Z) T , W  \right\rangle - \left\langle R(Z,T) V  ,  W  \right\rangle - \left\langle R(T,Z) W ,  V  \right\rangle - \left\langle R(Z,W) T ,  V  \right\rangle
}
{ =} { 2 \left\langle R(T,Z) V  ,  W  \right\rangle - \left\langle R(V,Z) T , W  \right\rangle - \left\langle R(Z,W) T ,  V  \right\rangle
}
}
{
\vergleichskettefortsetzungalign
{ =} { 2 \left\langle R(T,Z) V  ,  W  \right\rangle + \left\langle R(V,Z) W , T  \right\rangle + \left\langle R(Z,W) V  ,  T  \right\rangle
}
{ =} { 2 \left\langle R(T,Z) V  ,  W  \right\rangle - \left\langle R(W,V) Z , T  \right\rangle
}
{ =} { 2 \left\langle R(T,Z) V  ,  W  \right\rangle - \left\langle R(V,W) T , Z  \right\rangle
}
{ } {}
}{}{.}
Dies ergibt die Behauptung.
}

}

Aufgrund von
[[Riemannsche Mannigfaltigkeit/Levi-Civita-Zusammenhang/Krümmungsoperator/Eigenschaften/Fakt|Lemma 29.2  (2)]]
ist für Vektoren

\mavergleichskette
{\vergleichskette
{ v,w,z
}
{ \in }{ T_PM
}
{  }{
}
{    }{
}
{   }{
}
}
{}{}{}
in einem Punkt

\mavergleichskette
{\vergleichskette
{ P
}
{ \in }{ M
}
{  }{
}
{    }{
}
{   }{
}
}
{}{}{}
ein Vektor

\mavergleichskette
{\vergleichskette
{ R(v,w)(z)
}
{ \in }{ T_PM
}
{  }{
}
{    }{
}
{   }{
}
}
{}{}{}
wohldefiniert, da man ja die Vektoren durch differenzierbare Vektorfelder
\mathl{V,W,Z}{} mit

\mavergleichskette
{\vergleichskette
{ V(P)
}
{ = }{ v
}
{  }{
}
{    }{
}
{   }{
}
}
{}{}{}
etc.
\zusatzklammer {lokal} {} {}
realisieren kann und dann $R(V,W)(Z)$ im Punkt auswerten kann.

  \zwischenueberschrift{Die Schnittkrümmung}

\inputdefinition
{}
{

Es sei $M$ eine
\definitionsverweis {riemannsche Mannigfaltigkeit}{}{,}
die mit dem
\definitionsverweis {Levi-Civita-Zusammenhang}{}{}
auf dem
\definitionsverweis {Tangentialbündel}{}{}
$TM$ und dem zugehörigen
\definitionsverweis {Krümmungsoperator}{}{}
$R$ versehen sei. Dann nennt man zu
\definitionsverweis {linear unabhängigen}{}{}
Vektoren

\mavergleichskette
{\vergleichskette
{ v,w
}
{ \in }{ T_PM
}
{  }{
}
{    }{
}
{   }{
}
}
{}{}{}
über einem Punkt

\mavergleichskette
{\vergleichskette
{ P
}
{ \in }{ M
}
{  }{
}
{    }{
}
{   }{
}
}
{}{}{}
die Zahl

\mavergleichskettedisp
{\vergleichskette
{ K(v,w)
}
{ =} { { \frac{   \left\langle R(v,w)w , v  \right\rangle  }{  \left\langle  v  , v  \right\rangle  \left\langle  w  , w  \right\rangle  -\left\langle  v  , w  \right\rangle^2   } }
}
{ } {
}
{ } {
}
{ } {
}
}
{}{}{}
die
\definitionswort {Schnittkrümmung}{}
zu $v,w$ in $P$.

}

Diese Schnittkrümmung ist wohldefiniert, da im Fall der linearen Unabhängigkeit der Nenner nicht $0$ ist.

__NOEDITSECTION__

\inputfaktbeweis
{Riemannsche Mannigfaltigkeit/Levi-Civita-Zusammenhang/Schnittkrümmung/Tangentiale Ebene/Fakt}
{Lemma}
{}
{

\faktsituation {Es sei $M$ eine
\definitionsverweis {riemannsche Mannigfaltigkeit}{}{,}
die mit dem
\definitionsverweis {Levi-Civita-Zusammenhang}{}{}
auf dem
\definitionsverweis {Tangentialbündel}{}{}
$TM$ und dem zugehörigen
\definitionsverweis {Krümmungsoperator}{}{}
$R$ versehen sei.}
\faktfolgerung {Dann hängt die
\definitionsverweis {Schnittkrümmung}{}{}
zu zwei linear unabhängigen Vektoren

\mavergleichskette
{\vergleichskette
{ v,w
}
{ \in }{ T_PM
}
{  }{
}
{    }{
}
{   }{
}
}
{}{}{}
nur von der durch die Vektoren erzeugten Ebene

\mavergleichskettedisp
{\vergleichskette
{ \R v + \R w
}
{ \subseteq} { T_PM
}
{ } {
}
{ } {
}
{ } {
}
}
{}{}{}
ab.}
\faktzusatz {}
\faktzusatz {}

}
{

Wir betrachten den definierenden Ausdruck

\mavergleichskettedisp
{\vergleichskette
{ K(v,w)
}
{ =} { { \frac{   \left\langle R(v,w)w , v  \right\rangle  }{  \left\langle  v  , v  \right\rangle \left\langle  w  , w  \right\rangle  - \left\langle  v  , w  \right\rangle^2   } }
}
{ } {
}
{ } {
}
{ } {
}
}
{}{}{}
für die Schnittkrümmung. Der Ausdruck
\mathl{R(v,w)z}{} ist nach
[[Riemannsche Mannigfaltigkeit/Levi-Civita-Zusammenhang/Krümmungsoperator/Eigenschaften/Fakt|Lemma 29.2]]
linear in allen Komponenten. Wenn man $v$ oder $w$ mit einem Skalar

\mavergleichskette
{\vergleichskette
{ a
}
{ \neq }{ 0
}
{  }{
}
{    }{
}
{   }{
}
}
{}{}{}
multipliziert, so kann man sowohl im Zähler als auch im Nenner den Faktor $a^2$ rausziehen. Ferner gilt

\mavergleichskettealignhandlinks
{\vergleichskettealignhandlinks
{ \left\langle  R(v+w,w) w  , v +w   \right\rangle
}
{ =} { \left\langle  R(v,w)w +R(w,w)w  , v+w   \right\rangle
}
{ =} { \left\langle  R(v,w)w   , v   \right\rangle  +  \left\langle  R(w,w)w  ,  v   \right\rangle + \left\langle  R(v,w)w  ,  w   \right\rangle  +  \left\langle  R(w,w)w   , w   \right\rangle
}
{ } {
}
{ } {
}
}
{}
{}{.}
Nach
[[Riemannsche Mannigfaltigkeit/Levi-Civita-Zusammenhang/Krümmungsoperator/Eigenschaften/Fakt|Lemma 29.2  (4)]]
sind
\mathkor {} {\left\langle  R(w,w)w  ,  v   \right\rangle} {und} {\left\langle  R(w,w)w   , w   \right\rangle} {}
gleich $0$ und wegen
[[Riemannsche Mannigfaltigkeit/Levi-Civita-Zusammenhang/Krümmungsoperator/Eigenschaften/Fakt|Lemma 29.2  (5,6)]]
ist auch

\mavergleichskette
{\vergleichskette
{ \left\langle  R(v,w)w  ,  w   \right\rangle
}
{ = }{ 0
}
{  }{
}
{    }{
}
{   }{
}
}
{}{}{.}
Der Zähler der Schnittkrümmung ändert sich also nicht, wenn man $v$ durch $v+w$ ersetzt. Wegen

\mavergleichskettealigndrucklinks
{\vergleichskettealigndrucklinks
{ \left\langle v+w , v+w  \right\rangle \left\langle  w  , w  \right\rangle -\left\langle v+w , w  \right\rangle^2
}
{ =} { \left\langle  v  , v  \right\rangle  \left\langle  w  , w  \right\rangle+\left\langle  w  , w  \right\rangle \left\langle  w  , w  \right\rangle +2\left\langle  v  , w  \right\rangle \left\langle  w  , w  \right\rangle -\left\langle v  , w  \right\rangle^2-\left\langle w  , w  \right\rangle^2 -2 \left\langle v  , w  \right\rangle \left\langle w  , w  \right\rangle
}
{ =} { \left\langle  v  , v  \right\rangle \left\langle  w  , w  \right\rangle -\left\langle v  , w  \right\rangle^2
}
{ } {
}
{ } {
}
}
{}{}{}
ändert sich dabei auch nicht der Nenner.

}

\inputbemerkung
{}
{

Auf einer zweidimensionalen riemannschen Mannigfaltigkeit gibt es in jedem Tangentialraum nur den Raum selbst als zweidimensionalen Untervektorraum. Das bedeutet, dass im zweidimensionalen Fall
nach
[[Riemannsche Mannigfaltigkeit/Levi-Civita-Zusammenhang/Schnittkrümmung/Tangentiale Ebene/Fakt|Lemma 29.4]]
die Argumente in der
\definitionsverweis {Schnittkrümmung}{}{}
überflüssig sind. Deshalb fassen wir in dieser Situation die Schnittkrümmung als eine reellwertige Funktion auf der Mannigfaltigkeit auf.

}

__NOEDITSECTION__

\inputfaktbeweis
{Riemannsche Mannigfaltigkeit/Zweidimensional/Levi-Civita-Zusammenhang/Schnittkrümmung/Berechnung/Fakt}
{Lemma}
{}
{

\faktsituation {Für eine zweidimensionale
\definitionsverweis {riemannsche Mannigfaltigkeit}{}{}}
\faktfolgerung {ist die
\definitionsverweis {Schnittkrümmung}{}{}
auf jeder Karte gleich
\mathl{-  { \frac{   R^2_{121} }{ g_{11}  } }}{.}}
\faktzusatz {}
\faktzusatz {}

}
{

Wir können direkt die Situation auf einer Karte betrachten. Nach
[[Riemannsche Mannigfaltigkeit/Levi-Civita-Zusammenhang/Schnittkrümmung/Tangentiale Ebene/Fakt|Lemma 29.4]]
kann man im zweidimensionalen Fall die Schnittkrümmung mit einer beliebigen Basis ausrechnen, wir nehmen
\mathkor {} {\partial_1} {und} {\partial_2} {.}
Für die Schnittkrümmung ist unter Verwendung von
[[Riemannsche Mannigfaltigkeit/Levi-Civita-Zusammenhang/Krümmungsoperator/Eigenschaften/Fakt|Lemma 29.2  (5)]]

\mavergleichskettealign
{\vergleichskettealign
{ K( \partial_1, \partial_2)
}
{ =} { { \frac{   \left\langle R(\partial_1 ,\partial_2)\partial_2  ,  \partial_1   \right\rangle  }{  \left\langle  \partial_1  ,  \partial_1   \right\rangle  \left\langle  \partial_2  ,  \partial_2   \right\rangle  -\left\langle  \partial_1  ,  \partial_2   \right\rangle^2   } }
}
{ =} { { \frac{   \left\langle R(\partial_1 ,\partial_2)\partial_2  ,  \partial_1   \right\rangle  }{  g_{11} g_{22}  -g_{12}^2   } }
}
{ =} { -  { \frac{   \left\langle R(\partial_1 ,\partial_2)\partial_1  ,  \partial_2   \right\rangle  }{  g_{11} g_{22}  -g_{12}^2   } }
}
{ } {
}
}
{}
{}{}
Dabei ist gemäß
[[R^n/Offene Menge/Riemannsche Metrik/Levi-Civita-Zusammenhang/Krümmungsoperator/Berechnung/Fakt|Lemma 29.1]]

\mavergleichskettedisp
{\vergleichskette
{ R( \partial_1 , \partial_2)( \partial_1)
}
{ =} { R^1_{121} \partial_1+ R^2_{121} \partial_2
}
{ } {
}
{ } {
}
{ } {
}
}
{}{}{.}
Dies ergibt

\mavergleichskettealign
{\vergleichskettealign
{ K( \partial_1, \partial_2)
}
{ =} { - { \frac{   \left\langle R(\partial_1 ,\partial_2)\partial_1  ,  \partial_2   \right\rangle  }{  g_{11} g_{22}  -g_{12}^2   } }
}
{ =} { - { \frac{   R^1_{121} g_{12}+ R^2_{121} g_{22}  }{  g_{11} g_{22}  -g_{12}^2   } }
}
{ } {
}
{ } {
}
}
{}
{}{.}
Nach
[[Riemannsche Mannigfaltigkeit/Levi-Civita-Zusammenhang/Krümmungsoperator/Eigenschaften/Fakt|Lemma 29.2  (4)]]
ist

\mavergleichskettedisp
{\vergleichskette
{ \left\langle  R(\partial_1, \partial_1) \partial_1  ,  \partial_2   \right\rangle
}
{ =} { R^1_{121}  g_{11}  + R_{121}^2 g_{12}
}
{ =} { 0
}
{ } {
}
{ } {
}
}
{}{}{.}
Wir multiplizieren den Zähler der Schnittkrümmung mit $g_{11}$ und erhalten

\mavergleichskettealignhandlinks
{\vergleichskettealignhandlinks
{ g_{11}  \left(  R^1_{121} g_{12}+ R^2_{121} g_{22}  \right)
}
{ =} { g_{11}  \left(  R^1_{121} g_{12}+ R^2_{121} g_{22}  \right) - g_{12}  \left(  R^1_{121}  g_{11}  + R_{121}^2 g_{12}   \right)
}
{ =} { g_{11} g_{22} R^2_{121} - g_{12}^2 R^2_{121}
}
{ =} { \left(  g_{11} g_{22}  - g_{12}^2  \right) R^2_{121}
}
{ } {
}
}
{}
{}{.}
Also ist

\mavergleichskettedisp
{\vergleichskette
{ K( \partial_1, \partial_2)
}
{ =} { -  { \frac{   R^2_{121} }{ g_{11}  } }
}
{ } {
}
{ } {
}
{ } {
}
}
{}{}{.}

}

__NOEDITSECTION__

\inputfaktbeweis
{Fläche im Raum/Schnittkrümmung/Gaußkrümmung/Fakt}
{Lemma}
{}
{

\faktsituation {Es sei

\mavergleichskette
{\vergleichskette
{ W
}
{ \subseteq }{ \R^3
}
{  }{
}
{    }{
}
{   }{
}
}
{}{}{}
\definitionsverweis {offen}{}{}
und sei
\maabb {h} { W } { \R
} {}
eine zweifach
\definitionsverweis {stetig differenzierbare Funktion}{}{}
und

\mavergleichskette
{\vergleichskette
{ Y
}
{ = }{ h^{-1}(c)
}
{  }{
}
{    }{
}
{   }{
}
}
{}{}{}
die Faser zu

\mavergleichskette
{\vergleichskette
{ c
}
{ \in }{ \R
}
{  }{
}
{    }{
}
{   }{
}
}
{}{}{,}
wobei $h$ in jedem Punkt von $Y$
\definitionsverweis {regulär}{}{}
sei.}
\faktfolgerung {Dann ist die
\definitionsverweis {Gaußkrümmung}{}{}
von $Y$ gleich der
\definitionsverweis {Schnittkrümmung}{}{}
von $Y$.}
\faktzusatz {}
\faktzusatz {}

}
{

Nach
[[Riemannsche Mannigfaltigkeit/Zweidimensional/Levi-Civita-Zusammenhang/Schnittkrümmung/Berechnung/Fakt|Lemma 29.6]]
ist die Schnittkrümmung auf einer Karte gleich
\mathl{- { \frac{  R^2_{121}  }{ g_{11}  } }}{.} Wegen

\mavergleichskettedisphandlinks
{\vergleichskettedisphandlinks
{ R^2_{121}
}
{ =} { {\partial_1 } \Gamma^2_{12} - {\partial_2 } \Gamma^2_{11} +  \Gamma^1_{12} \Gamma_{1 1}^2 + \Gamma^2_{12} \Gamma_{1 2}^2 - \Gamma^1_{11}  \Gamma_{12 }^2  - \Gamma^2_{11} \Gamma_{2 2}^2
}
{ =} { -g_{11} K
}
{ } {
}
{ } {
}
}
{}{}{.}
ist

\mavergleichskettedisphandlinks
{\vergleichskettedisphandlinks
{ K
}
{ =} { { \frac{   1  }{ g_{11}  } }  \left(  \partial_2 \Gamma_{11}^2 - \partial_1 \Gamma_{12}^2 + \Gamma^1_{11} \Gamma^2_{12}  + \Gamma_{11}^2 \Gamma_{22}^2 - \Gamma_{12}^1 \Gamma_{11}^2-  \left(  \Gamma_{12}^2  \right)^2  \right)
}
{ } {
}
{ } {
}
{ } {
}
}
{}{}{}
was mit der Gaußkrümmung gemäß
[[Hyperfläche/Raum/Parametrisierung/Gaußkrümmung/Christoffelsymbole/Fakt|Satz 19.7]]
übereinstimmt.

}

\inputbeispiel{}
{

Für die
\definitionsverweis {euklidische Ebene}{}{}
$\R^2$ mit dem
\definitionsverweis {Standardskalarprodukt}{}{}
ist die
\definitionsverweis {Schnittkrümmung}{}{}
konstant gleich $0$.

}

\inputbeispiel{}
{

Für die
\definitionsverweis {Einheitskugel}{}{}

\mavergleichskette
{\vergleichskette
{ S^2
}
{ \subseteq }{ \R^3
}
{  }{
}
{    }{
}
{   }{
}
}
{}{}{}
mit der induzierten
\definitionsverweis {riemannschen Struktur}{}{}
ist die
\definitionsverweis {Schnittkrümmung}{}{}
konstant gleich $1$. Dies folgt aus
[[Kugeloberfläche/Gaußkrümmung/Beispiel|Beispiel 5.6]]
in Verbindung mit
[[Fläche im Raum/Schnittkrümmung/Gaußkrümmung/Fakt|Lemma 29.7]].

}

\inputbeispiel{}
{

\bild{ \begin{center}
\includegraphics[width=5.5cm]{ \bildeinlesungsvg {Parallel lines in Poincare's model of hyperbolic geometry} {svg} }
\end{center}
\bildtext {} }

\bildlizenz { Parallel lines in Poincare's model of hyperbolic geometry.svg } {} {Januszkaja} {Commons} {CC-by-sa 3.0} {}

Wir knüpfen an
[[Halbebene/Poincaré/Riemannsche Geometrie/Beispiel|Beispiel 18.8]]
und
[[Halbebene/Poincaré/Riemannsche Geometrie/Christoffelsymbole/Beispiel|Beispiel 26.7]]
an. Mit
[[R^n/Offene Menge/Riemannsche Metrik/Levi-Civita-Zusammenhang/Krümmungsoperator/Berechnung/Fakt|Lemma 29.1]]
ist

\mavergleichskettedisphandlinks
{\vergleichskettedisphandlinks
{ R^2_{121}
}
{ =} { {\partial_1 } \Gamma^2_{12} - {\partial_2 } \Gamma^2_{11} + \Gamma^1_{12} \Gamma_{1 1}^2 + \Gamma^2_{12} \Gamma_{1 2}^2 - \Gamma^1_{11} \Gamma_{12 }^2 - \Gamma^2_{11} \Gamma_{2 2}^2
}
{ =} { { \frac{   1  }{  y^2 } } - { \frac{   1  }{  y^2 } }+  { \frac{   1  }{  y^2 } }
}
{ =} { { \frac{   1  }{  y^2 } }
}
{ } {
}
}
{}{}{.}
Nach
[[Riemannsche Mannigfaltigkeit/Zweidimensional/Levi-Civita-Zusammenhang/Schnittkrümmung/Berechnung/Fakt|Lemma 29.6]]
gilt für die
\definitionsverweis {Schnittkrümmung}{}{}

\mavergleichskettedisp
{\vergleichskette
{ K
}
{ =} { - { \frac{   R^2_{121} }{ g_{11 } }}
}
{ =} { - { \frac{   { \frac{   1  }{  y^2 } }  }{  { \frac{   1  }{  y^2 } }  } }
}
{ =} { -1
}
{ } {
}
}
{}{}{,}
es liegt also konstante Schnittkrümmung $-1$ vor.

}

 [[Kategorie:Latexseite]]
